%! AP Statistics — Unit 2 Cover Sheet
\documentclass[10pt]{article}
\usepackage{apstats-article}
\usepackage{apstats-boxes}

%=============================================================================
\begin{document}

% ── Full-bleed navy banner ────────────────────────────────────────────────────
\begin{tikzpicture}[remember picture, overlay]
  \fill[navy] (current page.north west)
    rectangle ([yshift=-1.4in]current page.north east);
\end{tikzpicture}
\vspace{-1.3in}

\begin{center}
  {\color{white}\bfseries\fontsize{28}{34}\selectfont AP Statistics}\\[8pt]
  {\color{goldacc}\bfseries\Large Shepherd \quad 2026--2027}\\[14pt]
  {\color{white}\bfseries\LARGE Unit 2: Exploring Two-Variable Data}\\[4pt]
\end{center}

\vspace{0.30in}

% ── Unit overview ─────────────────────────────────────────────────────────────
\begin{tcolorbox}[colback=sky,colframe=navy,boxrule=0.9pt,arc=2mm,
    left=4mm,right=4mm,top=3mm,bottom=3mm]
  \textbf{Unit Overview.}\quad
  Unit 2 extends statistical thinking to the relationship between two variables.
  Students begin with two-variable categorical data, building two-way tables and
  comparing conditional distributions to assess association. They then shift to two
  quantitative variables, constructing scatterplots and describing association using
  direction, form, and strength. The unit develops correlation, simple linear
  regression, and residual analysis as tools for modeling and evaluating linear
  relationships, and closes with power and exponential transformations to linearize
  curved patterns.
\end{tcolorbox}

\vspace{0.22in}

% ── Lessons in this unit ──────────────────────────────────────────────────────
\renewcommand{\arraystretch}{1.6}
\begin{skillbox}[Lessons in This Unit]{goldbox}
\begin{tabularx}{\linewidth}{@{} c >{\bfseries}p{0.42\linewidth} X @{}}
  \textbf{\#} & \textbf{Title} & \textbf{Focus} \\
  \hline
  2.1 & Introducing Statistics: Are Variables Related?
    & Statistical questions about relationships; apparent vs.\ real associations. \\
  \hline
  2.2 & Representing Two Categorical Variables
    & Two-way tables; joint, marginal, and conditional relative frequencies. \\
  \hline
  2.3 & Statistics for Two Categorical Variables
    & Comparing conditional distributions; segmented bar charts; assessing association. \\
  \hline
  2.4 & Representing Two Quantitative Variables
    & Scatterplots; direction, form, and strength; explanatory vs.\ response variable. \\
  \hline
  2.5 & Correlation
    & Pearson's $r$; properties and interpretation; correlation $\ne$ causation. \\
  \hline
  2.6 & Linear Regression Models
    & Least-squares line $\hat{y}=a+bx$; writing and interpreting slope and intercept. \\
  \hline
  2.7 & Residuals
    & Calculating residuals; residual plots; assessing linearity. \\
  \hline
  2.8 & Least Squares Regression
    & Computing $b = r\,s_y/s_x$ and $a = \bar{y}-b\bar{x}$; $s$ and $r^2$; influential points. \\
  \hline
  2.9 & Analyzing Departures from Linearity
    & Power and exponential transformations; choosing and evaluating linearized models. \\
  \hline
\end{tabularx}
\end{skillbox}

\vspace{0.18in}

% ── Standards covered ─────────────────────────────────────────────────────────
\begin{skillbox}[AP Statistics Big Ideas \& Learning Objectives --- Unit 2]{greenbox}
\begin{tabularx}{\linewidth}{@{} >{\bfseries}p{0.14\linewidth} X @{}}
  VAR-1.D & Identify when an apparent association between two variables may be due
            to chance rather than a real relationship. \\
  UNC-1.P & Construct and interpret two-way tables; compute joint, marginal, and
            conditional relative frequencies to determine whether an association exists. \\
  DAT-1.I & Construct and interpret scatterplots for two quantitative variables;
            describe direction, form, and strength. \\
  DAT-1.J & Calculate and interpret the correlation coefficient $r$; understand its
            properties and limitations. \\
  DAT-1.K & Determine the least-squares regression line; interpret slope and
            $y$-intercept in context; use the line for prediction. \\
  DAT-1.L & Calculate and interpret residuals; use residual plots to assess whether a
            linear model is appropriate. \\
  DAT-1.M & Interpret the standard deviation of the residuals $s$ and the coefficient
            of determination $r^2$; identify outliers and influential points. \\
  DAT-1.N & Apply power or exponential transformations to achieve linearity; choose
            and evaluate the better model using residual plots and $r^2$. \\
\end{tabularx}
\end{skillbox}

\end{document}
